%==============================================================================
% APÊNDICE D — INSTALAÇÃO DO AMBIENTE
%==============================================================================
\chapter{Instalação e Configuração do Ambiente}
\label{apendice:instalacao}

Este apêndice guia a instalação do ambiente de desenvolvimento necessário
para executar todos os exemplos e projetos do livro.

\section{Opção 1: Google Colab (Recomendado para Iniciantes)}

O Google Colab oferece um ambiente Jupyter Notebook com GPU gratuita,
dispensando instalação local.

\begin{enumerate}
  \item Acesse \url{https://colab.research.google.com/}
  \item Faça login com sua conta Google
  \item Selecione \texttt{Arquivo $\rightarrow$ Abrir Notebook}
  \item Escolha um notebook do repositório do livro
  \item Ative a GPU em \texttt{Tempo de execução $\rightarrow$ Alterar
    tipo de tempo de execução $\rightarrow$ T4 GPU}
\end{enumerate}

\section{Opção 2: Ambiente Local (usuários avançados)}

\subsection{Instalação do Python}

\codigo{codigos/apendice-d-instalacao.sh}

\subsection{Criação do Ambiente Virtual}

Recomendamos o uso de um ambiente virtual para isolar as dependências:

\begin{lstlisting}[language=bash, numbers=none]
python3 -m venv odonto-ia-env
source odonto-ia-env/bin/activate  # Linux/Mac
# ou
odonto-ia-env\Scripts\activate.bat  # Windows
\end{lstlisting}

\subsection{Instalação das Dependências}

\begin{lstlisting}[language=bash, numbers=none]
pip install -r codigos/requirements.txt
\end{lstlisting}

\section{Verificação da Instalação}

Execute o seguinte script para confirmar que tudo está funcionando:

\lstinputlisting{codigos/cap01/test_pratica1_1.py}

\teste{Script de verificação executado com sucesso em ambiente local e Colab.}

\section{Comandos Úteis}

\begin{lstlisting}[language=bash, numbers=none]
# Verificar versão do Python
python3 --version

# Listar bibliotecas instaladas
pip list

# Atualizar pip
pip install --upgrade pip
\end{lstlisting}

\section{Solução de Problemas Comuns}

\begin{description}
  \item[Erro: TensorFlow não encontrado]
    Execute \texttt{pip install tensorflow}
  \item[Erro: PyTorch não encontrado]
    Acesse \url{https://pytorch.org/} para o comando de instalação adequado
  \item[Erro: OpenCV não encontrado]
    Execute \texttt{pip install opencv-python}
  \item[Erro de permissão no Linux]
    Use \texttt{pip install --user} ou um ambiente virtual
\end{description}